% Part: first-order-logic
% Chapter: syntax-and-semantics
% Section: covered-structures

\documentclass[../../../include/open-logic-section]{subfiles}

\begin{document}

\olfileid{fol}{syn}{cov}

\olsection{Gedekte \printtoken{P}{structure} voor eerste-ordetalen}

\begin{explain}
Bedenk dat een term \emph{gesloten} heet als hij geen !!{variable}s bevat.
\end{explain}

\begin{defn}[!!^{value} van gesloten termen]
Als $t$ een gesloten term van de taal~$\Lang L$ is en $\Struct M$ een
!!{structure} voor~$\Lang L$, wordt de \emph{!!{value}}~$\Value{t}{M}$
als volgt gedefinieerd:
\begin{enumerate}
\item Als $t$ enkel de !!{constant}~$c$ is, dan is $\Value{c}{M} = \Assign{c}{M}$.
\item Als $t$ de vorm $\Atom{f}{t_1, \ldots, t_n}$ heeft, dan
  \[
  \Value{t}{M} = \Assign{f}{M}(\Value{t_1}{M}, \ldots,
  \Value{t_n}{M}).
  \]
\end{enumerate}
\end{defn}

\begin{defn}[Gedekte !!{structure}]
Een !!{structure} heet \emph{gedekt} als elk element van het domein de
!!{value} van een gesloten term is.
\end{defn}

\begin{ex}
Zij~$\Lang L$ de taal met de !!{constant}s $\Obj{zero}$,
$\Obj{one}$, $\Obj{two}$, \dots, het binaire !!{predicate}~$<$ en de
binaire !!{function}s $+$ en $\times$. Neem als !!a{structure}~$\Struct
M$ voor~$\Lang L$ de structuur met domein $\Domain M = \{0, 1, 2, \ldots
\}$ en toekenningen $\Assign{\Obj{zero}}{M} = 0$,
$\Assign{\Obj{one}}{M} = 1$, $\Assign{\Obj{two}}{M} = 2$, enzovoort.
Voor het binaire relatiesymbool $<$ is de verzameling $\Assign{<}{M}$
de verzameling van alle paren $\tuple{c_1, c_2} \in \Domain{M}^2$
waarvoor $c_1$ kleiner is dan~$c_2$: zo geldt bijvoorbeeld
$\tuple{1, 3} \in \Assign{<}{M}$, maar $\tuple{2, 2} \notin
\Assign{<}{M}$. Voor de binaire !!{function} $+$ definiëren we
$\Assign{+}{M}$ op de gebruikelijke wijze---zo beeldt
$\Assign{+}{M}(2,3)$ af op~$5$---en analoog voor de binaire
!!{function}~$\times$. De !!{value} van $\Obj{four}$ is dus gewoon
$4$, en de !!{value} van $\times(\Obj{two},
+(\Obj{three},\Obj{zero}))$ (of, in infixnotatie, $\Obj{two} \times
(\Obj{three} + \Obj{zero})$) is
\begin{multline*}
\Value{\times(\Obj{two}, +(\Obj{three},\Obj{zero}))}{M} =\\
\begin{aligned}
& =\Assign{\times}{M}(\Value{\Obj{two}}{M}, \Value{+(\Obj{three}, \Obj{zero})}{M})\\
& = \Assign{\times}{M}(\Value{\Obj{two}}{M}, \Assign{+}{M}(\Value{\Obj{three}}{M},
\Value{\Obj{zero}}{M})) \\
& = \Assign{\times}{M}(\Assign{\Obj{two}}{M}, \Assign{+}{M}(\Assign{\Obj{three}}{M},
\Assign{\Obj{zero}}{M})) \\
& = \Assign{\times}{M}(2, \Assign{+}{M}(3, 0)) \\
& = \Assign{\times}{M}(2, 3) \\
& = 6
\end{aligned}
\end{multline*}
\end{ex}

\begin{prob}
Is $\Struct N$, het standaardmodel van de rekenkunde, gedekt? Licht uw
antwoord toe.
\end{prob}

\end{document}
